\chapter{Conclusões}
\label{cap:conclusions}

\section{Síntese do Trabalho Realizado}

Este projeto teve como objetivo central demonstrar que é possível gerir uma infraestrutura de serviços de rede de forma totalmente automatizada, reprodutível e segura, utilizando exclusivamente ferramentas \textit{open source}. O resultado final é uma plataforma robusta e funcional que corre numa máquina virtual Ubuntu Server 24.04 LTS, orquestrando um ecossistema completo de \textit{containers} Docker: um servidor DNS (Pi-hole), um \textit{reverse proxy} estrito (Nginx), uma VPN (WireGuard), uma nuvem privada corporativa (Nextcloud), um servidor de e-mail (Docker-Mailserver) e uma \textit{stack} avançada de observabilidade (Prometheus, Node Exporter, Loki, Promtail e Grafana).

Do ponto de vista técnico, o projeto atingiu todos os seus objetivos e superou as perspetivas iniciais. A infraestrutura é definida integralmente em código (ficheiros YAML do Docker Compose e \textit{Playbooks} Ansible), com segredos blindados pelo Ansible Vault e versionada no repositório do GitLab. A implantação automática a cada \textit{commit} pela \textit{pipeline} de CI/CD prova que qualquer falha catastrófica no servidor poderia ser recuperada em poucos minutos. Bastaria provisionar uma nova VM e executar o repositório, uma garantia de resiliência que, em infraestruturas geridas manualmente, simplesmente não existe.

\section{Aprendizagens e Contribuições}

O processo de implementação revelou que os maiores desafios técnicos raramente surgem da configuração dos serviços em si, mas sim das interações e restrições entre as diferentes camadas da arquitetura. Problemas como o \textit{bridging} incorreto do VMware (Secção \ref{sec:vmware_bridging}), os erros de indentação YAML no Netplan, o bloqueio do tráfego WireGuard pelo \textit{router} de fronteira, a partilha do \textit{socket} do Docker com o Promtail ou até a expansão em tempo real do volume LVM, ilustram bem esta realidade. Cada um destes incidentes exigiu um raciocínio estruturado de diagnóstico, eliminando hipóteses uma a uma até identificar a causa raiz.

A utilização do \texttt{tcpdump} para diagnosticar o problema de conectividade do WireGuard foi particularmente elucidativa: ao constatar que zero pacotes chegavam à interface de rede do servidor, foi possível descartar de imediato qualquer problema na configuração do \textit{container} ou do protocolo, focando o diagnóstico na barreira de rede externa (a configuração de DMZ no \textit{router}). Esta abordagem metódica --- de observar e analisar o tráfego antes de agir --- é um princípio fundamental da administração de sistemas que este projeto ajudou a consolidar na prática.

A integração do Ansible com a \textit{pipeline} GitLab CI/CD demonstrou como a filosofia de \textit{Infrastructure as Code} elimina a dicotomia histórica entre o que está documentado e o que está realmente instalado \cite{morris2016infrastructure}. Neste modelo, o repositório Git é, simultaneamente, a documentação viva e o estado real da infraestrutura.

\section{Comparação com Trabalhos Relacionados}

A abordagem adotada neste projeto partilha a filosofia de orquestração leve com outras soluções documentadas na literatura e na comunidade técnica. Em comparação com arquiteturas baseadas em Kubernetes \cite{docker_docs}, a solução Docker Compose apresenta uma complexidade operacional consideravelmente menor e é mais adequada para ambientes de pequena escala, como o laboratório académico ou PMEs. Para contextos com dezenas ou centenas de serviços e requisitos severos de \textit{auto-scaling}, o Kubernetes seria a escolha natural, mas introduz uma curva de aprendizagem e uma sobrecarga de gestão de \textit{hardware} que não se justificam neste cenário.

A escolha do WireGuard em detrimento do tradicional OpenVPN, fundamentada nas propriedades criptográficas modernas e na reduzida base de código apresentadas por Donenfeld \cite{donenfeld2017wireguard}, revelou-se extremamente acertada. A configuração foi significativamente mais simples, a integração com \textit{Dynamic DNS} fluiu sem atritos e o desempenho do túnel foi visivelmente superior em testes subjetivos de latência e navegação remota.

\section{Trabalho Futuro}

Com base na análise crítica dos resultados (Capítulo \ref{cap:test}) e considerando que os serviços principais e de segurança (Vault, Loki, E-mail e Nuvem) foram concluídos com sucesso, identificam-se as seguintes linhas de evolução natural para a plataforma:

\begin{itemize}
    \item \textbf{Cifragem HTTPS com Let's Encrypt:} Implementar a terminação TLS real no Nginx utilizando certificados públicos gratuitos gerados pelo \textit{Let's Encrypt} (via \texttt{certbot}). Em alternativa, migrar o encaminhamento web para o \textbf{Traefik}, que atua como \textit{reverse proxy} com renovação automática de certificados \cite{rfc8446}.

    \item \textbf{Alertas Automáticos Proativos:} Configurar regras de notificação no \textit{Alertmanager} do Grafana. O objetivo seria despoletar o envio automático de mensagens (por exemplo, via Telegram ou para o e-mail corporativo interno agora ativo) sempre que a utilização de CPU ultrapassasse os 85\% ou a disponibilidade de memória RAM caísse para valores críticos \cite{grafana_docs}.

    \item \textbf{Estratégia de Backups Off-site:} Uma vez que os dados do Nextcloud e as caixas de correio do Mailserver estão atualmente persistidos em volumes locais no host, a implementação de uma rotina automatizada de backup encriptado (utilizando ferramentas como o Restic ou o BorgBackup) para um armazenamento externo protegeria a infraestrutura contra falhas físicas do disco da máquina virtual.
\end{itemize}

\section{Considerações Finais}

Este projeto demonstrou que a conceção e a gestão profissional de infraestruturas de rede já não são um exclusivo de grandes organizações com largas equipas de TI. As ferramentas \textit{open source} disponíveis atualmente --- Docker, Ansible, GitLab, Prometheus, WireGuard --- alcançaram um nível de maturidade e documentação que permite a um único administrador de sistemas, munido dos conhecimentos técnicos adequados, construir e operar uma plataforma de serviços robusta, cifrada e totalmente automatizada. 

Fica provado que o principal obstáculo na transição para um \textit{data center} moderno não é de cariz tecnológico, mas sim metodológico: passa por adotar desde o primeiro dia a disciplina inegociável do versionamento contínuo, da configuração declarativa e da documentação sistemática dos problemas encontrados.